Fix \(\mathcal S^{\mathrm{cal}}\in\mathfrak M_{\mathrm{cal}}\) and write
\(q_h=P(W_h=1)\).  The bounds in \(\Theta_{\mathrm{cal}}\) ensure that every
threshold in \(\mathrm{def{:}calibrated\text{-}bridge\text{-}family}\) lies in
\([0,1]\).  Independence of \(U_W,U_M,U_Y\) from one another and from the
remaining exogenous variables gives, at every conditional event of positive
probability,
\[
 \mathbb E[W\mid H=h,A=a,M=m]=q_h,
 \qquad
 \mathbb E[M\mid H=h,A=1]=c+dq_h.
 \tag{5}
\]
Consistency and the treatment-zero outcome equation therefore imply
\[
 \begin{aligned}
 \mathbb E[Y\mid H=h,M=m,A=0]
 &=\alpha+\tau m+\beta q_h\\
 &=\mathbb E[h_0(W,m)\mid H=h,M=m,A=0].
 \end{aligned}
 \tag{6}
\]
This calculation actually holds at every positive \((H,M,A=0)\) cell and hence,
in particular, at every target-relevant cell guaranteed positive by
\(\mathrm{ass{:}latent\text{-}a0\text{-}transport\text{-}support}\).  Similarly,
conditional independence of \(W\) and \(M\) given \((H,A=1)\) and (5) give
\[
 \begin{aligned}
 \mathbb E[h_0(W,M)\mid H=h,A=1]
 &=\alpha+\tau(c+dq_h)+\beta q_h\\
 &=\alpha+\tau c+(\beta+\tau d)q_h\\
 &=\mathbb E[h_1(W)\mid H=h,A=1].
 \end{aligned}
 \tag{7}
\]
The target-relevant conditioning events in (7) are positive by
\(\mathrm{ass{:}latent\text{-}a1\text{-}transport\text{-}support}\).

The same common-shock calculation, now conditional on \(H_0=h\), gives
\[
 \mathbb E[Y_{0,h,M_{1,h},W_h}\mid H_0=h]
 =\alpha+\tau(c+dq_h)+\beta q_h
 =\mathbb E[h_1(W_h)\mid H_0=h].
\]
The maintained treatment independence makes \(H_0\perp A\), and consistency
identifies \((H_0,W_{H_0})\) with \((H,W)\) in arm zero.  Summing over the three
positive hidden states proves
\[
 \psi(\mathcal S^{\mathrm{cal}})=\mathbb E[h_1(W)\mid A=0].
 \tag{8}
\]

We next verify the observed equations rather than assuming them.  Conditional on
\((H,M,A=0)\), equation (6) holds for both sides.  The exclusions
\(Z\mathbin{\perp\!\!\!\perp}(Y,M)\mid(H,A)\) and
\(W\mathbin{\perp\!\!\!\perp}(A,Z,M)\mid H\) allow conditioning additionally on
\(Z=z\) without changing the two within-\(H\) expectations.  Averaging over
\(H\mid Z=z,M=m,A=0\) proves the first four equations of
\(B(P)x=b(P)\).  In the treatment-one arm, the same exclusions make the joint
conditional expectation in (7) invariant to \(Z\) within each hidden state.
Averaging (7) over \(H\mid Z=z,A=1\) proves the last two equations.  The observed
denominator assumptions make all six conditionals well defined.  Thus the vector
induced by \((h_0,h_1)\) lies in both
\(\mathcal H_{\mathrm{lat}}(\mathcal S^{\mathrm{cal}})\) and
\(\mathcal H_{\mathrm{obs}}(P_{\mathcal S^{\mathrm{cal}}}^O)\).

All member properties of \(\mathfrak M_{\mathrm{WBC}}^+\), except the intersection
property just proved, are already member properties of
\(\mathfrak M_{\mathrm{cal}}\).  Hence
\(\mathfrak M_{\mathrm{cal}}\subseteq\mathfrak M_{\mathrm{WBC}}^+\), and
\(\mathrm{lem{:}intersection\text{-}implies\text{-}solvability}\) gives observed
solvability.  If \(\ell(P)\in\operatorname{row}(B(P))\), equation (8),
\(\mathrm{lem{:}latent\text{-}bridge\text{-}identifies}\), and
\(\mathrm{lem{:}null\text{-}annihilator}\) imply for every observed solution
\[
 \psi(\mathcal S^{\mathrm{cal}})=\ell(P)^\top x=r(P)^\top b(P).
\]
The membership and rank assertions for \(\mathcal S^+\) are exactly the proved
content of \(\mathrm{prop{:}positive\text{-}rank\text{-}two\text{-}success}\).

Finally, \(\mathrm{lem{:}proximal\text{-}id\text{-}scope}\) records that the cited
Proximal ID theorem retains injectivity and produces a recursive interventional
margin, while \(\mathrm{lem{:}proximal\text{-}path\text{-}specific\text{-}scope}\)
records that the cited path-specific result uses an observed recanting witness, a
distinct latent baseline confounder, and three completeness conditions.  The
present derivation instead concerns only the displayed scalar same-unit target on
\(\mathcal G_{\mathrm{WBC}}\), and it constructs neither comparator's graph-level
algorithm or inference theory.  Conversely, the positive rank-two member just
proved violates the injective proxy premise used by those routes.  These opposing
domain and conclusion differences establish the two stated non-subsumption
comparisons.